\section{逻辑门按真值表改变输出}

逻辑门接收一个或多个 0/1 输入，并按固定规则产生输出。AND 门只有在两个
输入都为 1 时输出 1；NOT 门把 0 变成 1、把 1 变成 0。四种输入组合可以
完整写成真值表：

\begin{osgridtabular}{L{0.16\textwidth}L{0.16\textwidth}L{0.22\textwidth}}
\(A\) & \(B\) & \(A\land B\)\\ \hline
0 & 0 & 0\\ \hline
0 & 1 & 0\\ \hline
1 & 0 & 0\\ \hline
1 & 1 & 1\\ \hline
\end{osgridtabular}

真值表列出所有输入，因此不会漏掉边界组合。后面读标志位、权限位和状态机时，
同样要把所有可能输入列清楚。

\section{组合逻辑不记得过去}

若输出只取决于此刻输入，这类电路叫组合逻辑。两个 1-bit 数相加时，XOR
给出和位，AND 给出进位：
\[
  S=A\oplus B,\qquad C=A\land B.
\]
多路选择器则根据选择位，从几组输入中挑一组送到输出。CPU 的寄存器选择、
ALU 操作数选择和设备地址译码都会使用这类结构。

组合逻辑需要传播时间。输入刚改变时，内部路径尚未全部稳定，输出可能短暂
出现过渡值。第一次主线只在时钟边沿到来前等待它稳定；hazard、建立时间和
亚稳态的定量分析放在数字逻辑深化附录。

\section{反馈让电路保存一个状态}

如果把门电路的输出反馈到输入，过去的输出就可能影响现在。锁存器利用反馈
保存一个 bit：写入动作结束后，输入即使恢复，内部状态仍能保持。这里第一次
出现“状态”——只看当前外部输入已不足以预测输出，还必须知道此前保存的值。

状态必须有初始值。复位信号把锁存器、计数器或状态机带到规定状态，否则上电
后可能从任意组合开始。CPU 的复位向量、内核对象的初始枚举和磁盘格式的 magic
都在解决同类问题：让后续步骤拥有共同起点。

\section{时钟决定什么时候更新寄存器}

D 触发器在规定的时钟边沿采样输入，并把结果保持到下一次采样。把许多触发器
并排放置，就得到能够一次保存多个 bit 的寄存器。8-bit 寄存器保存一个字节，
64-bit 寄存器保存 64 个 bit；位宽增加不会改变基本动作。

计数器可以在每个时钟边沿把寄存器加一。程序计数器也是寄存器，但它保存下一
条指令的位置。顺序执行时它增加指令长度，跳转时它被替换成目标地址。

\section{寄存器和逻辑门可以组成小型 CPU}

最小教学 CPU 需要程序计数器、指令寄存器、通用寄存器、ALU 和控制器。
程序计数器给出地址，存储器返回指令字节，控制器决定读取哪些寄存器，ALU
完成加法或逻辑运算，结果再写回寄存器。

我们暂时把这个过程写成五步：
\[
  \text{取指}\rightarrow\text{译码}\rightarrow\text{执行}
  \rightarrow\text{访存}\rightarrow\text{提交}.
\]
真实 x86 会把多条指令重叠执行，但软件最终观察到的寄存器、内存和异常仍要
满足架构规则。下一章先沿一条不重叠的教学路径，看 CPU 怎样真正取得第一条
指令。
